\begin{summarybox}
\textbf{\textcolor{CSummary}{一句话核心}}

\textbf{同曲线两点相减，平方差的“和”给中点、“差”给斜率：}这是点差法的全部发动机。

\medskip
\textbf{\textcolor{CSummary}{最短复现（不依赖死记）}}
\[
\begin{array}{rcl}
\text{椭圆：}&\dfrac{x_0}{a^2}+\dfrac{y_0}{b^2}k=0
&\Longrightarrow\ k=-\dfrac{b^2x_0}{a^2y_0},\\[6pt]
\text{双曲线：}&\dfrac{x_0}{a^2}-\dfrac{y_0}{b^2}k=0
&\Longrightarrow\ k=\dfrac{b^2x_0}{a^2y_0},\\[6pt]
\text{抛物线：}&y_0k=p
&\Longrightarrow\ k=\dfrac{p}{y_0},\\[6pt]
\text{圆：}&x_0+y_0k=0
&\Longrightarrow\ k=-\dfrac{x_0}{y_0}.
\end{array}
\]
以上斜率公式均须 $y_0\ne0$；默认确实存在两个不同实交点。

\medskip
\textbf{\textcolor{CSummary}{两步考场检查}}
\begin{itemize}[leftmargin=*,itemsep=0.4em,topsep=0.3em]
\item \textbf{第一步（抓信号）：}出现“弦中点、对称、求弦方程、求中点轨迹” $\Rightarrow$ 先点差求候选方向。
\item \textbf{第二步（看逆推）：}若题目已告知弦存在，直接计算；若给中点\emph{反求}弦，必须确认“两不同实交点”或检查有效范围。
\end{itemize}

\medskip
\textbf{\textcolor{CSummary}{最容易忘记的两个反例}}

\textbf{反例一：} $M=O$ 的椭圆中心弦可以水平、竖直或倾斜，不能把 $0/0$ 当作“竖直”。
\textbf{反例二：} 双曲线 $x^2/4-y^2/5=1$ 的 $M(1,1)$ 可以算出候选斜率 $5/4$，但不存在相应实弦。

\medskip
\textbf{\textcolor{CSummary}{证明选用顺序}}

\textbf{最快：}两式相减；\quad
\textbf{门槛最低：}两端点写成 $M\pm(u,v)$；\quad
\textbf{常规：}联立直线后用韦达；\quad
\textbf{进阶统一：} $\nabla F(M)\cdot(u,v)=0$。
\end{summarybox}
